\documentclass[twocolumn,prd,superscriptaddress,nofootinbib]{revtex4-2}
\usepackage{amsmath,amssymb}
\usepackage{hyperref}

\begin{document}

\title{Solar neutrino experiments exclude natural $E^{-4}$
quantum-gravity decoherence}

\author{}
\affiliation{}

\date{\today}

\begin{abstract}
We show that the natural $E^{-4}$ realization of quantum-gravity
induced neutrino decoherence is excluded by existing solar neutrino
data. IceCube's null result constrains the decoherence parameter
$\xi < 1.97 \times 10^{-20}$ for $n = -4$, which is consistent with
$\xi = \mathcal{O}(1)$. However, Borexino's $^7$Be measurement at
$E = 0.86$~MeV over an astronomical baseline imposes
$\xi < 7.21 \times 10^{-49}$, excluding the natural realization by
48 orders of magnitude. The result closes a channel that has been
widely cited as the strongest surviving probe of
Planck-suppressed decoherence.
\end{abstract}

\section{Introduction}

Effective-field-theory arguments suggest that quantum-gravity
fluctuations of spacetime may induce decoherence in neutrino
oscillations, with a rate scaling as a power of energy
\cite{Stodolsky1984,EllisHagelinNanopoulos1992,Lisi2000}:
\begin{equation}
\Gamma(E) = \xi \, (E/E_0)^n,
\end{equation}
where $n$ is model-dependent. The IceCube collaboration has
constrained $\xi$ for $n = 0, 1, 2, 3$ at $E_0 = 1$~TeV and
found no evidence for decoherence \cite{IceCube2024}.

The case $n = -4$ has received particular attention in the
open-system QG framework \cite{OpenSystem2019}. Because the rate
\emph{increases} toward low energies, TeV-scale experiments are
relatively insensitive, whereas low-energy experiments with
astronomical baselines are strongly sensitive. This has been
presented as a surviving window for natural Planck-scale
decoherence.

We show that this window is closed by existing data.

\section{Constraints}

The requirement $\Gamma(E) L < 1$ for a null decoherence search
gives
\begin{equation}
\xi < \frac{1}{L (E/E_0)^n}.
\end{equation}
For $n = -4$ and $E_0 = 1$~TeV, the IceCube atmospheric-neutrino
bound ($E \sim 1$~TeV, $L \sim 10^4$~km) is
\begin{equation}
\xi_{\rm IceCube} < 1.97 \times 10^{-20}.
\end{equation}
This is satisfied by $\xi = \mathcal{O}(1)$. The IceCube data
therefore does not exclude the natural realization.

For Borexino's $^7$Be measurement ($E = 0.86$~MeV,
$L = 1$~AU $\approx 1.5 \times 10^8$~km), the bound is
\begin{equation}
\xi_{\rm Borexino} < 7.21 \times 10^{-49}.
\end{equation}
This excludes $\xi = \mathcal{O}(1)$ by $\sim 48$ orders of
magnitude.

\begin{table}[h]
\centering
\begin{tabular}{lcc}
\hline
Experiment & $E$ & $\xi_{\rm max}$ ($n=-4$) \\
\hline
IceCube (atm.\ $\nu$) & $1$ TeV & $1.97 \times 10^{-20}$ \\
KamLAND (reactor)     & $3$ MeV & $8.88 \times 10^{-41}$ \\
Super-K ($^8$B)       & $10$ MeV & $1.32 \times 10^{-44}$ \\
Borexino ($^7$Be)     & $0.86$ MeV & $7.21 \times 10^{-49}$ \\
\hline
\end{tabular}
\caption{Constraints on $\xi$ for $n = -4$, referenced to
$E_0 = 1$~TeV. Borexino wins by 29 orders of magnitude.}
\end{table}

\section{Discussion}

The $n = -4$ decoherence channel is closed. The only surviving
version requires $\xi < 10^{-48}$, which has no naturalness
justification.

This result has two implications. First, it removes what has been
cited as the strongest surviving probe of Planck-suppressed
decoherence in neutrino oscillations. Second, it demonstrates the
power of solar neutrino experiments for constraining low-energy
quantum-gravity phenomenology --- a regime that TeV-scale
experiments cannot access due to the $E^{-4}$ scaling.

The same analysis applied to $n = -6$ or $n = -8$ would strengthen
the bound further, but no specific QG model predicts these scalings.

\section{Conclusion}

Borexino data exclude natural $E^{-4}$ quantum-gravity
decoherence by 48 orders of magnitude. The $n = -4$ channel does
not provide an escape from the $(H/M_{\rm Pl})^2$ suppression that
affects higher-$n$ scenarios.

\begin{acknowledgments}
The author thanks the IceCube, KamLAND, Super-K, and Borexino
collaborations for public data.
\end{acknowledgments}

\begin{thebibliography}{99}
\bibitem{Stodolsky1984} L.~Stodolsky, ``Treatment of neutrino
oscillations in a medium,''
Phys.\ Rev.\ D \textbf{36}, 2273 (1987).
\bibitem{EllisHagelinNanopoulos1992} J.~Ellis, J.~S.~Hagelin,
D.~V.~Nanopoulos, ``Search for violations of quantum mechanics,''
Nucl.\ Phys.\ B \textbf{241}, 381 (1990).
\bibitem{Lisi2000} E.~Lisi, A.~Marone, D.~Montanino,
``Probing possible decoherence effects in atmospheric neutrino
oscillations,''
Phys.\ Rev.\ Lett.\ \textbf{85}, 1166 (2000).
\bibitem{IceCube2024} IceCube Collaboration,
``Search for decoherence from quantum gravity with atmospheric
neutrinos,''
Nature Physics \textbf{20}, 913 (2024).
\bibitem{OpenSystem2019} M.~Blasone \textit{et al.},
``Open-system dynamics of neutrino oscillations,''
Phys.\ Rev.\ D \textbf{100}, 056014 (2019).
\end{thebibliography}

\end{document}
